
\section{Appendix}

%\subsection{Initial Prompt Used for Design Generation}
%\label{sec:appendix-prompt}

This appendix provides the initial natural language prompts used for generating the evaluated accelerator workloads presented in Section~\ref{04_Results&Evaluation}.


%******************************************************************************************************

\subsection{Element-wise Vector Multiplication}
\emph{
I would like to create a hardware accelerator design. The accelerator should be able to take two input vectors: X and Y, both of length L.
The accelerator should perform an element-wise multiplication
operation and produce an output vector Z.
The accelerator has two AXI-Stream based interfaces for loading
X and Y data into custom X and Y buffers.
The accelerator should also have a fixed length parameter L.
Once the data is loaded, the accelerator should execute the
element-wise multiplication in parallel and store the results
in buffer Z within the compute module.
The loading should be performed using a load module. Finally, the results should be written back to main memory using a store module that outputs via an AXI-Stream interface.
Create the accelerator description using SystemC and SECDA. The compute module should be capable of performing L operations in parallel.
}

\subsection{2D Convolution}
\emph{
I would like to create a hardware accelerator design.
The accelerator should be able to take two input matrix: one for weights and another for inputs.
The weights have the dimensions in follow data order: input channels (IC), output channels (OC), kernel width (KW) and then kernel height (KH) .
The input matrix has the dimensions in follow data order: input channels (IC), input width (IW) and then input height (IH).
The output matrix has the dimensions in follow data order: output channels (OC), output width (OW) and then output height (OH).
We assume padding is 0 and stride is 1 and dilation is 1 for the convolution operation.
So the output width (OW) and output height (OH) can be calculated as follows:
OW = IW - KW + 1
OH = IH - KH + 1
The accelerator should perform convolution operation and produce a output matrix Z.
The accelerator has two AXI-Stream based interfaces for loading weights and inputs into custom weights and inputs buffers.
All dimension are fixed and defined as parameters in the accelerator design.
Once the data is loaded, the accelerator should be abe to execute the CONV operation in parrallel and save it in buffer Z, within the compute module.
The loading should be done with load module.
And finally the storing of the results back to main memory should be done in the store module which writes to an output AXI-Stream.
Create the accelerator description using SystemC and SECDA.
}

\subsection{Matrix Transpose}
\emph{
I would like to create a hardware accelerator design.
The accelerator should be able to take a input matrix X of dimensions (M, N), transpose it into an output matrix Z of dimensions (N, M).
Where M, N are not fixed and are provided as data at runtime, before matrix X is loaded into the accelerator. The accelerator should be able to handle any dimensions of M and N, as long as they fit within the memory constraints (buffer size) of the accelerator design.
The accelerator has an AXI-Stream based interface for loading the input matrix X into a custom input buffer.
Once the data is loaded, the accelerator should be able to execute the transpose operation and save the result in buffer Z, within the compute module.
The loading should be done with load module.
The Compute module should be able to perform the transpose operation. It should d
And finally the storing of the results back to main memory should be done in the store module which writes to an output AXI-Stream..
Create the accelerator description using SystemC and SECDA.
}